% Related lecture: SC2005-L05
% Source: Part 3 (CPU Scheduling), pp. 11--13, FCFS lecture example
% Human verified: Yes

\exercisemeta
  {SC2005-L05-EX}
  {2026-08-27}
  {Part 3 (CPU Scheduling), pp.~11--13，FCFS lecture example}
  {题干、Gantt charts 与答案均已核对}
  {已确认（2026-08-27）}

\exercisequestion{SC2005-L05-E01}{FCFS Order 与 Convoy Effect}

\textbf{对应知识点：}\par
\relatedtopic{SC2005-L05-T03}{Scheduling Objectives 与 Time Metrics}；
\relatedtopic{SC2005-L05-T04}{FCFS 与 Convoy Effect}

\subsubsection*{题目（Lecture Example）}

三个 processes 均在 $t=0$ 到达，CPU burst lengths 为
\[
  B_1=24,\qquad B_2=3,\qquad B_3=3.
\]
分别按 ready-queue order（就绪队列顺序）$P_1,P_2,P_3$ 与 $P_2,P_3,P_1$ 使用 FCFS。画出 Gantt charts，计算各 process 的 waiting time 与 average waiting time，并解释两种结果的差异。

\begin{checkedsolution}
两种 schedule 为：
\begin{center}
  \centering
  \resizebox{0.96\textwidth}{!}{\begin{tikzpicture}[
  x=0.34cm,y=0.9cm,
  every node/.style={font=\small},
  long/.style={draw=ntuRed,very thick,fill=red!7,minimum height=0.8cm,align=center},
  short/.style={draw=noteBlue,very thick,fill=blue!9,minimum height=0.8cm,align=center},
  tick/.style={font=\scriptsize,anchor=north}
]
  \node[anchor=east,font=\bfseries] at (-0.6,1.5) {$P_1,P_2,P_3$};
  \node[long,minimum width=24*0.34cm] at (12,1.5) {$P_1$};
  \node[short,minimum width=3*0.34cm] at (25.5,1.5) {$P_2$};
  \node[short,minimum width=3*0.34cm] at (28.5,1.5) {$P_3$};
  \foreach \x in {0,24,27,30} {\draw (\x,1.02) -- (\x,0.84); \node[tick] at (\x,0.82) {\x};}
  \node[anchor=west,text=ntuRed] at (31.0,1.5) {$\overline W=17$};

  \node[anchor=east,font=\bfseries] at (-0.6,-0.2) {$P_2,P_3,P_1$};
  \node[short,minimum width=3*0.34cm] at (1.5,-0.2) {$P_2$};
  \node[short,minimum width=3*0.34cm] at (4.5,-0.2) {$P_3$};
  \node[long,minimum width=24*0.34cm] at (18,-0.2) {$P_1$};
  \foreach \x in {0,3,6,30} {\draw (\x,-0.68) -- (\x,-0.86); \node[tick] at (\x,-0.88) {\x};}
  \node[anchor=west,text=noteBlue] at (31.0,-0.2) {$\overline W=3$};
\end{tikzpicture}
}

  \smallskip
  \textit{相同 workload、不同 ready-queue order 下的 FCFS schedules。}
\end{center}

顺序 $P_1,P_2,P_3$：
\[
  (W_1,W_2,W_3)=(0,24,27),\qquad
  \overline W=\frac{0+24+27}{3}=17.
\]
因为 arrival times 都为 $0$，completion times 也就是 turnaround times：
\[
  (T_1,T_2,T_3)=(24,27,30),\qquad \overline T=27.
\]

顺序 $P_2,P_3,P_1$：
\[
  (W_1,W_2,W_3)=(6,0,3),\qquad
  \overline W=\frac{6+0+3}{3}=3,
\]
\[
  (T_1,T_2,T_3)=(30,3,6),\qquad \overline T=13.
\]
本题为 single-burst、nonpreemptive scheduling，所以每个 process 的 response time 等于其 waiting time。

第一种顺序中，长任务 $P_1$ 迫使两个短任务等待，形成 convoy effect。第二种顺序只是改变题目给定的入队顺序，并不表示 FCFS 会主动按 burst length 重排；主动优先短任务属于下一讲的 SJF。
\end{checkedsolution}

\textbf{检查点：}先根据 arrival/order 画 Gantt chart，再由每个 process 的实际 start time 计算 $W_i=S_i-A_i$；不能先按 burst length 排序后仍称为 FCFS。
